% results.tex

\section*{Results}

\subsection*{A target-centered framework for using simulation labels in Tm prediction}

The target task was nanobody melting-temperature (Tm) regression from
amino-acid sequence. Computational labels were treated as auxiliary prediction
tasks, not as stand-ins for experimental Tm. Each model used the pretrained
protein sequence model ESM2 to turn a sequence into a learned representation
that a Tm head and an auxiliary head shared, so the auxiliary task could reshape
the representation while the final objective stayed anchored to experimental Tm
(Fig.~\ref{fig:learning_framework}).

We used the benchmark split as a deliberately low-data target setting, with 57
experimental Tm labels for training, 114 for validation, and 396 for held-out
testing. To keep the comparison fair across very different auxiliary labels,
each source was tuned independently: for every source and every encoder regime we
searched the model architecture, the auxiliary-head coupling, and the learning
rate, dropout, and weight decay, selecting on the experimental Tm validation set
alone. Only the selected configuration was then evaluated on the held-out Tm test
set, and improvements were measured as paired differences in absolute error on the
same test examples. This target-centered protocol guards against a trap specific
to simulation-to-experiment learning: an auxiliary task can be learnable and
physically sensible yet do nothing for the experimental phenotype. Under this
protocol the tuned Tm-only baselines were 7.23$^\circ$C MAE with a frozen encoder
and 6.55$^\circ$C MAE with a fine-tuned (hot) encoder. Two design axes organize the
results: whether a simulation label transfers at all is set by its \emph{data
design} (Fig.~\ref{fig:data_design}), and how deeply it transfers is set by its
\emph{physical observable} (Fig.~\ref{fig:physical_observable}).

\subsection*{Data design: matching the sequence scan makes an MD label transfer}

With each source tuned to the same Tm endpoint, two all-atom labels clearly
transferred: the FEP mutation free energy and --- new to this work --- an MD
native-contact stability label computed on the \emph{same} 1MEL and 4IDL mutation
scans used for FEP. The MD native-contact label uses one physical observable in
two different data designs, which isolates the effect of data design from the
choice of quantity. Computed over an earlier screen of diverse nanobodies, the
label did not improve Tm: its held-out error sat at or slightly above the
corresponding Tm-only reference (Fig.~\ref{fig:data_design}a, upper group).
Computed instead over the FEP-matched mutation scans --- the same scaffolds, all of
constant length --- the identical observable improved Tm and matched FEP with a
frozen encoder (Fig.~\ref{fig:data_design}a, lower group). Because the two groups
come from different training series with different absolute baselines, each is
shown against its own Tm-only reference and the two are not cross-compared
numerically.

The difference is a data-design confound, not the label arithmetic
(Fig.~\ref{fig:data_design}b). The diverse screen spans sequence lengths from 58 to
461 residues and carries a nonzero label--length correlation ($r = +0.13$), a
length signal the auxiliary task can latch onto instead of the stability signal;
the matched scans are of constant length, so no such confound exists. Referencing
each mutant's native-contact value to its wild type (a $\Delta Q$) versus using the
raw value is immaterial here, because a constant offset is absorbed by the min--max
scaling applied to every label. With the confound removed, both all-atom labels
scale below the Tm-only baseline as the number of auxiliary labels grows: with a
frozen encoder the FEP curve stays below the baseline throughout (7.13 at 20
labels to 7.01 at 320), and the matched-scan MD curve falls from above it to below
it, crossing by 320 labels (7.32 to 7.03; Fig.~\ref{fig:data_design}c). At 320
labels the matched MD label improved held-out Tm by $-0.20^\circ$C
(90\% CI $-0.34$ to $-0.05$), statistically indistinguishable from FEP in the same
regime ($-0.22^\circ$C; paired FEP-minus-MD difference $-0.03^\circ$C, 90\% CI
$-0.21$ to $+0.16$). Matching the auxiliary source's sequence design to the target
neighborhood, not the specific dynamical quantity, is what lets the MD label
transfer.

\begin{figure*}[t]
\centering
\includegraphics[width=\textwidth]{fig_outline02_data_design}
\caption{\textbf{Data-design axis: a matched mutation scan makes the MD label transfer.}
(a) Held-out test MAE for the MD native-contact label computed over a diverse
nanobody screen (upper) versus the FEP-matched 1MEL/4IDL mutation scan (lower),
each with its own Tm-only and FEP references. The two groups are from different
training series (different absolute baselines) and are not compared numerically
across groups; the point is null-versus-positive for the same observable.
(b) The confound made explicit: MD native-contact value versus sequence length.
The diverse screen has a broad length spread (58--461 residues) and a nonzero
label--length correlation ($r = +0.13$); the matched scans are constant length.
(c) Frozen-encoder label-count sweeps for FEP and the FEP-matched MD native-contact
label; the dashed line is the Tm-only frozen baseline. Both labels fall below the
baseline as labels increase, and the matched MD label crosses it by 320 labels.
Points and intervals show ensemble MAE and 90\% bootstrap intervals.}
\label{fig:data_design}
\end{figure*}

\subsection*{Physical observable: free energy transfers deeper than native contacts}

Tuned to the same endpoint, the sources formed a clear frozen-encoder hierarchy
(Fig.~\ref{fig:physical_observable}a): FEP (7.01) and the matched-scan MD label
(7.03) led, both improving on the Tm-only baseline (7.23); ThermoMPNN was a weak
third (7.09, 90\% CI crossing zero); and the Rosetta-family scores --- the plain
Rosetta mutation score and Rosetta scores on ESM2-proposed or random variant sets
--- sat at the baseline (7.23, 7.31, and 7.22). But a matched data design was
necessary, not sufficient: the two all-atom labels differed in how deeply they
transferred across encoder regimes (Fig.~\ref{fig:physical_observable}b). The FEP
free energy improved Tm whether the encoder was frozen or fine-tuned (frozen
$-0.22^\circ$C, 90\% CI $-0.36$ to $-0.08$; hot $-0.15^\circ$C, 90\% CI $-0.30$ to
$-0.01$), reaching the best held-out accuracy overall (6.40$^\circ$C MAE, hot). The
MD native-contact label improved Tm only with a frozen encoder; once the encoder
was fine-tuned it gave no net gain and, at small label counts, degraded prediction
(paired $+0.47^\circ$C at 20 labels, 90\% CI $+0.29$ to $+0.65$), recovering only
to the baseline by 320 labels. Directly, FEP beat the MD label in the hot regime
($-0.18^\circ$C paired, 90\% CI $-0.36$ to $-0.01$) while tying it when frozen.

The natural reading is that the alchemical free energy carries enough information
about the mutation's effect on stability to beneficially reshape the sequence
encoder, whereas a native-contact persistence label can usefully condition a fixed
representation but cannot steer encoder fine-tuning toward the Tm phenotype
(Fig.~\ref{fig:physical_observable}c). The physical content of the label sets the
\emph{depth} of transfer. The Rosetta-family scores fit the weaker side of this
picture: they did not robustly improve Tm in either regime. In particular the two
Rosetta-scored variant sets (ESM2-proposed and random variants) are ordinary
comparators; under tuning the ESM2-proposed set did not outperform the random set,
so we make no claim about physics-scored sequence proposals as a design loop.
Enlarging the ESM2 encoder from 8M to 35M or 650M parameters did not by itself
reproduce the FEP gain (Supplementary Fig.~4), so model capacity does not
substitute for a matched physical label.

\begin{figure*}[t]
\centering
\includegraphics[width=\textwidth]{fig_outline03_physical_observable}
\caption{\textbf{Physical-observable axis: depth of transfer set by the simulated quantity.}
(a) Tuned frozen-encoder held-out test MAE for every source; the dashed line is the
Tm-only baseline. FEP and the matched-scan MD label lead; the Rosetta-family scores
sit at the baseline.
(b) Paired $\Delta$MAE versus Tm-only for FEP, the MD native-contact label, Rosetta,
and ThermoMPNN under frozen versus hot encoders; values left of zero improve on
Tm-only. FEP crosses below zero in both regimes; the MD native-contact label only
when frozen; Rosetta and ThermoMPNN sit near zero.
(c) Interpretation: the free-energy label reshapes the shared encoder (deep;
transfers frozen and hot), whereas the native-contact label only refines the fixed
trunk (shallow; helps frozen, not hot). Points and intervals show ensemble MAE and
90\% bootstrap intervals.}
\label{fig:physical_observable}
\end{figure*}

Across every label we tested, the improvements were selective rather than
additive. Two axes organize the pattern: a matched data design decides whether a
simulation label transfers at all, and the physical observable decides how deep
that transfer goes. Under both, only the FEP mutation free energy (in both
regimes) and the FEP-matched MD native-contact label (with a frozen encoder)
robustly improved low-data nanobody Tm prediction.
